% =====================================================================
%  R. Nielsen, ON SELF-UNDERSTANDING IN EVIL, WITH REFERENCE TO SHAKESPEARE'S RICHARD THE THIRD
%  In: For Idea and Reality.  A Journal, edited by R. Nielsen, Bj. Bjørnson and Rud. Schmidt.  Volume for 1869.
%  Copenhagen: Chr. Steen & Søn; printed by Louis Klein.  pp. 521-554.  (Om Selvforstaaelse i det Onde, in For Ide og Virkelighed, Aargang 1869.)
%  ENGLISH TRANSLATION, from transcription.tex (Danish).
%  GENERATED by ebook-method/for-idee-1869/tr_build.py --write from this template and the translators' parts in
%  ebook-method/for-idee-1869/tr/ (TRANSLATOR-BRIEF.md; terms: GLOSSARY.md, the titles and four key terms approved by
%  Hans 2026-10-13).  Conventions: TRANSLATION-PLAYBOOK.md.  Page numbers of the 1869 print in the margin.
% =====================================================================
\documentclass[12pt,a4paper,openany]{book}
\usepackage[utf8]{inputenc}
\usepackage[T1]{fontenc}
\usepackage{libertinus}
\usepackage{libertinust1math}
\usepackage[english]{babel}
\usepackage[protrusion=true,expansion=true]{microtype}
\usepackage{setspace}
\usepackage[margin=1.3in]{geometry}
\usepackage{fancyhdr}
\usepackage{hyperref}
\hypersetup{pdftitle={On Self-Understanding in Evil}, pdfauthor={Rasmus Nielsen}, hidelinks}
\setlength{\headheight}{15pt}
\pagestyle{fancy}
\fancyhf{}
\fancyhead[LE,RO]{\thepage}
\fancyhead[C]{\textit{R. Nielsen: On Self-Understanding in Evil (1869)}}
\setstretch{1.3}
\usepackage{marginnote}
\renewcommand{\thefootnote}{\arabic{footnote})}   % the print numbers each page's note ¹) (p. 553: *))
\newcommand{\opage}[1]{\marginnote{\footnotesize\textit{[#1]}}}
\newcommand{\essayhead}[2]{\begin{center}{\large #1}\ifx&#2&\else\\[0.3em]{\normalsize #2}\fi\\[0.5em]\rule{3em}{0.4pt}\end{center}\medskip\noindent}
\newcommand{\signature}[1]{\par\smallskip\hfill #1\hspace*{2em}\par}

\begin{document}

% the volume's title page (PDF 7)
\begin{titlepage}
\begin{center}
\vspace*{3em}
{\Huge For Idea and Reality.}\\[1em]
\rule{8em}{0.4pt}\\[1em]
{\large A Journal,}\\[0.8em]
{\small edited}\\[0.4em]
{\footnotesize by}\\[0.6em]
{\large R. Nielsen,}\\[0.4em]
{\large Bj. Bjørnson and Rud. Schmidt.}\\[1em]
\rule{8em}{0.4pt}\\[1em]
{\large Volume for 1869.}\\[1em]
\rule{8em}{0.4pt}\\[4em]
Copenhagen.\\[0.4em]
Chr. Steen \& Søn, Publishers.\\[1.5em]
{\small Printed by Louis Klein.}\\[3em]
{\small Translated from the Danish}
\end{center}
\end{titlepage}

\essayhead{On Self-Understanding in Evil,}{with Reference to Shakespeare's Richard the Third.}
% --- p. 521 ---
\setcounter{footnote}{0}%
\opage{521}When a bird of prey swoops down upon a sea creature, and the contending
forces on both sides are about equal,
the outcome of the struggle will depend chiefly on the element
in which it is fought. If the inhabitant of the sea succeeds
in drawing its winged enemy down beneath the waves, then
it wins a double advantage: it fights in its own element,
and the element fights for it. Now it is not merely in
animal life that the strength of the forces in a struggle is conditioned
by the element; in the life of the spirit, too, something corresponding shows itself.
The spiritual forces in faith and knowledge are, with regard
to principles and presuppositions, representations and concepts,
so different that the corresponding ideas must
be said to move in two absolutely dissimilar regions,
and the struggle between them must every time depend on the element
in which it is conducted. To draw the religious thoughts of life
down into the stream of knowledge is arbitrarily to take sides for
the antireligious; for knowledge fights in its own objective element,
and the element fights for it. It is an age-old
error, a thousand-year-old theological deception, to confound
the elements, to deify the objective, and to drown
the eagle of faith in the waters of science.

% --- p. 522 ---
\setcounter{footnote}{0}%
\opage{522}It is a principle which anyone who wishes can
understand, that a matter of conscience is religious and moral,
but not scientific; or is there perhaps some reasonable
human being who in real life confounds a
concern of conscience with a scientific problem?
Nor, surely, the other way round! When a Jesper Ridefoged
feels in his conscience that the earth is not round,
one laughs; and why? Because one involuntarily feels
that physics and conscience must each speak its own language, each have
its own circle of thought with its own foundation of its own
grounds. That scientific questions cannot immediately
be submitted to the verdict of conscience
means that science itself is impersonal, that the element of
knowledge is objective. The religious and ethical truths
are, in the last instance as in the first, truths of conscience;
all true and thorough understanding of
these truths rests on self-understanding, but the element of
self-understanding is subjectivity. Hereby the
objective in religion and morality is by no means weakened,
still less denied; only this is to be emphasized, that
the judgment upon it is subjective, for the judge is in the human being's
self. If one fails to hear this inner judge, or
reduces him to a subordinate official and seeks
the supreme judge in the scientific, then it is easy to
foresee what the consequence must be. It is thus of great importance,
not only for the sake of religiosity but also for that of cultural history and humanity,
to give oneself an account
of the difference between scientific truth and truth of conscience.

In all countries where science and literature keep
pace with the times, there now exist literary documents which
% --- p. 523 ---
\setcounter{footnote}{0}%
\opage{523}clearly prove that theology, with its learned baggage,
is without principle, that its medieval role is played out.
But, one will perhaps object, if this
is so, how then is one to explain
that a fact so manifest, so
deeply affecting culture and humanity, has not yet really dawned
on the cultivated public at large? The matter is easily explained.
By forcing faith down into the objective element of knowledge,
theology has from generation to generation bewitched itself and its age
into the fixed idea that theology and religion are inseparable;
``if theology falls, religion falls''---that is the
dogmatic watchword meant to strike into the civilized world
terror of atheism and respect for theology. How
false this watchword is can be seen from the dispute about faith and
knowledge. The supreme principle of theology and of atheism
is, namely, one and the same. ``If theology falls,
religion falls''---only on this presupposition can it
be explained that an atheistic criticism which has scientifically
dissolved theology can imagine that it has at the
same time dissolved positive religion and thus
won a double victory. That theology's ``objective faith'' has
fallen to the weapons of knowledge is manifest to all who know, and
yet, strangely enough, a public secret. How
such a secret can maintain itself in an enlightened
age, in spite of so many telescopes and so many clear
eyes, is readily understood when one considers that it is not
only in the well-understood interest of the theologians but also of the atheists
that the common people, ``the great
mob'', do not get to peep too deep into the pot. From
both sides one finds one's account in covering the fact
with a thin veil. As to theology,
% --- p. 524 ---
\setcounter{footnote}{0}%
\opage{524}it would indeed be a suicide if it trumpeted forth its
own defeat and had it made known by advertisements
that it really is a sham science; it becomes
the business of its opponents to see to the announcement. One
would now think that among these opponents
the atheists, with their double victory, must stand in the foremost
rank. So it was, too, in the last century;
but in the present one the situation has strangely
changed. A psychological experience has been made.
The atheist, when he duly reflects, will see that he
cannot be served by the theological science of faith
losing all credit in the judgment of the people, for when the
enemy he lives by combating is dead and gone,
he has essentially nothing more to do; his esoteric
knowledge is a negative power that dare not expose its positive
foundation. As long as atheism has the metaphysical
veil over it, it is not unlike the goddess of Sais;
one feels a shudder before the dark secret, the
enigmatic, petrifying truth. But when the veil is lifted,
the likeness is over. The youth does not
go mad with horror; on the contrary, he becomes peevish,
because there is no horror, because his expectation
is disappointed: ``Is this the truth itself, is this the goddess's
countenance? Such a face one can meet with at any
dance-hall''! The flat denial can content itself
with the testimony of the five senses. To understand
the doctrine that there is no personal God and no
immortality of the soul, the practical consciousness of culture
need not trouble itself with obscure philosophical phrases;
such a wisdom of the commonplace one
breathes in with the common air. From atheism one may
% --- p. 525 ---
\setcounter{footnote}{0}%
\opage{525}therefore expect a complete unveiling just as little as from
theology. Under such circumstances public
opinion, whether it at a given time
inclines to theology or to atheism or perhaps---according as the display is in the ``shop''---to both
at once, must feel within itself that it is without insight,
without conviction, only an opinion, and so the secret
can well be preserved.

The public secret has, however, one good
about it that we ought not to fail to recognize. There runs through
the civilized world a sense that the dispute
for and against the religious dogmas is impractical, that
the practical side of the matter, that on which it really turns,
is morality and moral self-understanding; for,
it is said, ``he whose conduct is pure and right, his faith
cannot possibly be bad''. Theories that attack revealed
religion---call themselves Deism, Pantheism,
Atheism, or whatever they please---are by no means
dangerous to society, so long as they do not attack morality.
From the standpoint of the state and of humanity,
a capable, sound and vigorous moral doctrine is far to be preferred
to obscure, mystical speculations about revelation. Let a sober philosophical thinker, it is said,
furnish us a defensible
ethics, a practical moral doctrine that maintains duty and
right and gets to the bottom of the opposition of good and evil,
and he shall be welcome to us: we shall certainly not
trouble him with pressing questions about confession of faith.
That the view which in the popular sense can so
easily be stamped as atheism does not just, as regards the moral,
always stand behind orthodox doctrinal systems,
the example of the elder Fichte has
% --- p. 526 ---
\setcounter{footnote}{0}%
\opage{526}sufficiently shown. This great, moral, popularly minded thinker
was, when one could not otherwise get
at him, finally accused before the government as an atheist.
That helped, speaking in the civil sense, for he was
as Professor of Philosophy at the University of Jena actually
deposed from his office; but what it signified in the scientific
and moral sense can be seen from his manly,
open defense. Because a human being gives up faith
in a personal God, he does not therefore give up faith
in a divine world order. ``This faith'', says
Fichte, ``one must not imagine as an arbitrary
assumption, which a human being can make or not, as
it pleases him, as a free resolution to hold for
true what the heart wishes because it wishes it, as
a completion or substitute for convincing
grounds by calling hope to one's aid. What is grounded
in reason is simply and plainly necessary; and what is not
necessary is for that very reason contrary to reason. To
assume such a thing is imagination and dream, however piously
one may dream. The true good is cognized only under
the condition that I find myself free from all influences
of the world of the senses, absolutely active in myself,
thus as a power exalted above everything sensible. This
freedom is not indeterminate; it has its goal, only that it
does not receive it from without, but sets it for itself
by itself. Of this freedom and this its vocation
I cannot doubt without giving up myself. I cannot
doubt, I say---cannot even conceive the possibility
that it should not be so, that that
inner voice could deceive, that it would first have to be authorized
and grounded from elsewhere. Here no
reasoning avails, no reinterpreting or artificial interpretation.
% --- p. 527 ---
\setcounter{footnote}{0}%
\opage{527}That voice is absolutely positive and categorical. I cannot
get further without disturbing my own inner being;
I cannot go further precisely because I cannot
will to go further. Here lies that which sets its limit to the
otherwise untamed flight of ratiocination, that which
binds the spirit, that which binds the heart; here is the point
which unites thinking and willing into One and brings harmony
into my being. This is the true faith; the moral
order is the \emph{Divine} which we assume.
It is constructed by acting rightly. This is the only
possible confession of faith: freely and impartially to accomplish
what duty each time commands, without doubting and reasoning
over the consequences. Thereby the Divine becomes
living and real to us, each of our actions is accomplished
in the presupposition of this, and the consequences are preserved therein''.

One sees from this utterance where Fichte is heading,
and in what his moral consciousness is grounded.
If Fichte were only in agreement with himself about the
element in which the Idea of the good is to move and live, then
we could safely follow his guidance; for that he
earnestly wishes to maintain faith in the good as that which is in itself
true and divine, there is no doubt. But
precisely in the confusion of the spheres of faith and knowledge,
that is, in the confounding of the elements, lies his self-contradiction.
He rejects faith in a personal God;
this is by no means out of defiance, unbelief or ungodliness,
it is by virtue of scientific grounds, it is because a
personal God, according to his concept, must be a limited
God and a limited God is unthinkable. We
hear him in his defense: ``You attribute personality
and consciousness to God. What do you then understand by personality and
consciousness? Surely that which you have found in yourselves
% --- p. 528 ---
\setcounter{footnote}{0}%
\opage{528}and designated by this name? But that you in no
way think or can think this without limitation
and finitude, the slightest attention
to your construction of this concept can
teach you. By attributing to this being that predicate
you make it, then, a finite being, a being that is
like yourselves, and you have not, as you wished, thought God, but
only multiplied yourselves in thought''. But
if we now ask our philosopher on what ground we are then to
believe in the reality of the good and in a divine
government of the world, then it seems as if he suddenly
wished to move us over into a wholly different sphere, another
element. ``The conviction of our moral vocation'',
he says, ``itself already proceeds from a moral
disposition, and is faith, and one says so far quite
rightly: the element of all certainty is faith.---Thus
it had to be; for morality, however certain it is,
can simply only be constructed by itself, by no means
by any logical compulsion of thought''. Here one sees the inconsistency:
the question of the personality of God is to be decided
in science, about this, then, one may well reason;
but the question of the validity of morality
is not to be decided in science, for about this one may
not reason. The inconsistency becomes greater still when
we note that the author, having established
the principle of the self subjectively in faith---thus outside
science---nonetheless wants its development
from ``the simply and plainly necessary'' carried over into science.
When, then, someone says: I cannot give up
morality without giving up myself, he is declared
free of the summons of science; here no
logical compulsion of thought applies; but if someone says: I cannot
% --- p. 529 ---
\setcounter{footnote}{0}%
\opage{529}give up the personal God without giving up myself,
then science ought to summon him, for the personal God
is deposed by ``logical compulsion of thought''.

Only moral zeal to carry through an unconditionally
rational self-understanding in the good could have misled
a thinker like Fichte into committing so glaring
an inconsistency. The ethical error that is here at hand
can best be illuminated if, instead of speaking in general
of the general good and considering whether there
is a rational self-understanding in the good, we cast a glance at
the opposite of the good and ask whether there is not also
a rational self-understanding in evil. So long as the talk
is merely of the good, one easily settles into repose with general
considerations. The good, it is said, is One
with the true nature of man; in choosing the good
I essentially choose my own self; that I do the good
for the sake of the good means that I, in the deepest sense,
do it for my own sake, etc. Since it goes so
smoothly with such fine, general phrases, one
is easily tempted to treat truths of conscience
as scientific doctrines and to draw the highest
concerns of personality down into an impersonal
element. Otherwise when the talk is of evil. Here
the pure general considerations will not suffice. If
the essence of evil is to be comprehended in science, then the explanation
in the end amounts to this, that evil is only a deception,
an appearance. To understand oneself in evil would,
according to this way of viewing it, have to be regarded as just as impossible
as to do evil because it is evil. Only by
exchanging the general considerations from a distance for a real
view from near, that is: for a conception answering to idea and reality
of the personal in character
% --- p. 530 ---
\setcounter{footnote}{0}%
\opage{530}and will, will it be possible for thought to reach the root
of evil. A character portrait illuminating from this point of view
we have in Shakespeare's Richard the Third,
the knight of radical evil, a royal Satan of
flesh and blood, an incarnate devil.

What is peculiar to Richard's character is not that
he is bloodthirsty, cruel and commits crimes; in
this respect he is only a son of his time. The
time in which Richard lives is an evil time: violence and insolence,
murder and perjury belong to the order of the day. The medieval
state with its strong, self-willed aristocracy
is in the fifteenth century given over to passions of party
that threaten it with ruin; the struggle between
the white and the red rose, a civil war of more than thirty years,
in which threescore persons of the royal house and half the English nobility lost
their lives, is a bloody chapter in England's history. No, that
Richard is a tyrant is not the characteristic thing; the
characteristic lies in the consciousness with which he
exercises his tyranny. The acting persons in his surroundings,
whatever misdeeds they commit, have yet
a kind of excuse in their excited passions;
whirled into the vengeful war of the moment, all are more or
less blinded, only not Richard. It is a trait
of his terrible superiority that whereas none among
his fellow combatants, friend or foe, is really able to see through
him, he sees through them all, looks down upon
them, despises them and uses them as blind tools;
yes, what is more---for his demonic gaze is just as inward as outward---in seeing through
the others he sees through himself as well
% --- p. 531 ---
\setcounter{footnote}{0}%
\opage{531}and has, as it seems, a complete self-understanding in
evil.

The monologue with which the first scene of the first act begins
lets us gain a deep insight into Richard's soul.
The tumult of war is his element, he is bored in peace;
the effeminate court of his brother Edward the Fourth is
loathsome to him. Not from moral indignation, but because
he has his own plans, and then because, ``ugly-
stamped'', treated like a stepmother by Nature, he has no luck
with the ladies of the court, he declares that he will be a
villain ``and hate the idle pleasures of these days''.---``Plots
and conspiracies I have laid by dark omens,
dreams and libels, which shall set my brother Clarence and
King Edward at hatred and enmity against one another,
and if King Edward keeps his word, as surely as I
am cunning, false and faithless, Clarence shall this very day be
shut up in prison.'' So we hear at once from the beginning what
Richard is, and what he wants: he is a resolute man,
a man who is at one with himself. The horrifying
purpose with which Richard appears, and the manner
in which he motivates it, has occasioned a misunderstanding
which we must here reject. His words have been understood
as if in his natural ugliness he sought an
excuse for the malice with which he will avenge
himself on human society, from which as a
deformed being he is excluded. So A. W. Schlegel.
``Malice'', says this author, ``is nothing other than
a deliberately unconscientious egoism; yet
it can never quite lay aside the form of morality, because
this is law for all thinking beings; it must seek
to trace its perverse way of acting back to principles.
% --- p. 532 ---
\setcounter{footnote}{0}%
\opage{532}Although Richard knows very well the abomination of his disposition
and his mission from hell, he nevertheless seeks
to justify himself to himself by specious grounds; since
the happiness of being loved is denied him, what is left to him
but the happiness of ruling?'' This remark,
which is incidentally not without a certain psychological
fineness, does not fit Richard; it fits
scorned love, but he who has never loved
anyone but himself alone cannot feel the torments of a scorned
love. The love of which Richard
speaks is also not of such a nature that it can be regarded
as a virtue. When he says that he cannot ``play the lover'',
that he cannot ``court himself in the mirror'' and
``is not made to strut before a wanton nymph'',
he by no means means that it is anything virtuous
or good that is lost to him. Richard does not at all think
of a distinction between virtue and vice, but rather of a choosing between different
vices; he is not fitted for the one vice, and
therefore he chooses the other;---``since I cannot play the lover
and entertain these fair ladies, I am determined
to be a villain''; it is in the form of an excuse,
as if one were to say: ``Since I cannot steal,
I am determined to be a robber.'' It is literally the case
that Richard knows what a villain is, and \emph{in spite of that},
or rather \emph{on that very account}, has determined to be a
villain. ``So he wills evil, then, because it is
evil, what a monstrous contradiction!'' Certainly it is
an inconceivable contradiction so long as we regard him from
the general objective standpoint of knowledge, but by no means
when we understand him personally. So long as the word
% --- p. 533 ---
\setcounter{footnote}{0}%
\opage{533}villain is a designation for the evil man in strict opposition
to the good man, in other words: so long as the difference
between good and evil is not yet effaced from
the ground up, so long it will certainly be impossible
for a human being to will evil because it is evil.
But the difference between good and evil Richard has
personally effaced. By effacing this difference he has
by no means come to stand on neutral ground
as one who is neither good nor evil; on the contrary, he has
become sunk in evil; for to deny the difference
between good and evil is to devote oneself to evil. How
then did Richard originally come to abolish
the separation of good and evil? By denying
God, thus by taking an atheistic standpoint! But Richard's atheism has a peculiar coherence;
it shuns the objective element and keeps
consistently to the subjective. Richard's grounds are
not scientific, but unconditionally personal. To
write a critical metaphysical piece in which he abolished
revealed religion and announced a new era
when the human race was to be made blessed by ``an immanent
development of its immanent principle'', he had indeed
no time; to introduce a new liturgy with the formula: there
is no God, and Richard is his prophet, he did not
consider expedient, so he took the matter subjectively,
so subjectively that he contented himself with ``becoming
transparent to himself'', without becoming transparent to others.
The inner certainty with which Richard has given up God
shines forth precisely from the insolence with which he feigns
fear of God and uses religion, where it can serve
him as a means, for the attainment of his end. His
% --- p. 534 ---
\setcounter{footnote}{0}%
\opage{534}thoughts on this are expressed plainly enough in the monologue (First
Act, third scene, toward the end): ``I do the wrong, and first begin the brawl; the secret evil
I spread abroad, its burden I lay on others' shoulders\dots{}.
Now they believe it, and constantly whet me
to revenge on Rivers, Vaughan and on Grey. Then I sigh
and show with a scripture text that God bids us repay evil
with good. Thus I hide the naked
villain with old odd ends stolen out of the Bible, and
seem a saint when most I play the devil''. To what degree
his hypocrisy can shine we see then especially from the divine
comedy he plays with the citizens of London.
(Third Act, seventh scene.) Clarence is murdered; Edward
is dead, and his young sons kept in the Tower,
Rivers, Grey and Vaughan and even the noble, unsuspecting
Hastings have all mounted the scaffold; thus
Richard has prepared everything, and now it only matters that
his accession to the throne may get the appearance of occurring at
the people's invitation and be as legitimate as possible.
Richard must of course pretend that he is most unwilling to
take the heavy burden of government upon himself. Seized by dark
forebodings, the citizens then also are unwilling to choose him.
But Buckingham, Richard's copy and tool, knows how to
build the future on ``a holy plan''. ``Play'', he says to
Richard, ``the timorous; let yourself be long entreated before
you are seen. Stand, my lord, between two clergymen and take a
prayer book in your hand''\dots{} When now the citizens with their
Lord Mayor and Aldermen come, they are of course
turned away by the pious Richard, absorbed in godly contemplation.
So that the good faithful souls may nevertheless not
lose heart altogether because the ``noble Duke''
% --- p. 535 ---
\setcounter{footnote}{0}%
\opage{535}hesitates to receive them, Buckingham now explains the
real connection of the matter, saying: ``No, my lord,
this prince is no Edward; he does not stretch himself
on the soft day-bed, but on his knees in godly contemplation.
He toys not with a pair of wanton wenches, but
with two clergymen he searches deep.'' When then the Lord
Mayor at last sees Richard on a balcony between two
bishops, all doubt vanishes, and he exclaims in pious
rapture: ``See yonder his Grace between two bishops,
two pillars of virtue for a Christian prince, who
guard against the glitter of vanity! and see, a prayer book
is in his hand, the true ornament by which one
knows piety'' \dots{} To become what he afterwards himself
calls, ``the Lord's anointed'', Richard may well feel
a need to search with two clergymen, and, if need
be, to go to confession to two bishops: a man who daily
confesses so sincerely to himself may well have a few
confessions to spare for the bishops. He is a perfect
hypocrite; hypocrisy is a kind of denial of God, perfect
hypocrisy a thoroughgoing denial of God.

Therefore: Richard is an atheist; from this standpoint
his self-understanding and the consequences in his character
must be judged. If it now appears that an atheistic
ethics is not in a position to master Richard, then it is
thereby proved at once that atheistic ethics does not
possess the principle from which a moral opposition between
the good and evil can be developed and maintained. We
shall now try how the matter stands.

When God is denied, thus the will of God excluded,
man, seen in the Idea, is unconditionally his own
master, his own lawgiver, his own judge. This
% --- p. 536 ---
\setcounter{footnote}{0}%
\opage{536}must of course not be understood as if the individual in his
inner being---for as regards the outer, no
explanation is needed here---could arbitrarily make himself
whatever he will: the individual will is conditioned by and so far
bound to the individual nature, and the development of character
is at once a personal development of both
nature and will. It avails nothing in this connection to
take shelter under the metaphysical universal and talk at
length of the universally human, the universal nature,
the universal reason, the universal thought, the universal will;
for in every real human being the universal is limited
and assumes a peculiar shape; therefore
every human individuality has indeed its nature, its reason,
its thought and its will: the universal is not excluded,
but it is individualized. We cannot now
say that the real human nature, nature existing
in real human beings, is either unconditionally good
or unconditionally evil; it is, as is well known, a mixture
of both. As we must assume that in the evil nature,
however dark and corrupt it may be, there must yet always
be something good, since an absolutely evil nature would
be un-nature, and an absolute un-nature an impossibility, so
we must likewise assume that in the good nature,
however noble and pure it may otherwise be, in so far as
it is of Adamic principle, there must be something evil.
But then it is expressly admitted hereby that the evil
in human nature has its universality just as well as
the good. But when the universal can just as well
correspond to evil as to the good, then it is in vain
to want to maintain the moral demand by referring
the individual to what is universal in human nature.
% --- p. 537 ---
\setcounter{footnote}{0}%
\opage{537}The difference between the good and the evil nature is relative
and has as its extremes on the one side such
individuals in whom the good has an original preponderance
over the evil, on the other such as in
whom the reverse is the case; this difference acquires
a far-reaching significance when the task is to determine
the moral principle. Let us hear Fichte again, the man
of whom Goethe, though his fine worldly wisdom could not
approve the ethicist's energetic conduct, himself admits
that he was ``one of the most capable personalities
that has ever been seen'', a man in whose ``way of thinking
there is in a higher respect nothing to be criticized'': so
we feel that Fichte speaks out of the innermost of his soul when
he declares that he cannot doubt moral
freedom and its vocation without giving up himself,
and thus points to ``the point which unites thought
and will into One and brings harmony into his being.''
He wills the good as he wills his own true will; he
cannot go further without disturbing his own
inner being; he has his self-understanding in the good, because
his nature is good. If we turn next to Richard,
we get the same explanation with the same ground
and yet in a wholly opposite direction. Richard, too, has
the law for his actions unconditionally in himself; he too
wills what he, according to his peculiar natural endowments,
must will. As little as a Fichte can change
his nature, just as little can a Richard change his;
if he tries, he sinks and perishes instantly. From this point of view the self-contemplation
in which he with evident satisfaction emphasizes
the ugliness of his shape---``Of man's image
% --- p. 538 ---
\setcounter{footnote}{0}%
\opage{538}am an abridgment; deceitfully the hypocritical nature
described my features. I was sent into the world before my time,
kicked out of my mother's womb, scarcely half finished, so lame,
so unfinished, that dogs bark when I linger by them''---has no
other meaning than to serve as an empty excuse, because he does not
want what men call the good. Richard does not complain, does not
excuse himself; he takes himself as he is, he does not wish to
be otherwise or to be another: he is ``himself
alone.'' His outward ugliness is a symbol of a
devilish inwardness; but he is by no means dissatisfied with his
exterior; he has no ground for that, quite the contrary! The deformity
is for the demonic Richard what the repulsive
exterior, the snub nose, the protruding eyes, the warts,
the hanging belly and the clumsy feet were for the
ironic Socrates,---an excellent mask that suits
his role. The question here is not what is
ugly and what is beautiful, any more than what is
right and what is wrong: might is right. Richard's right
is the might of self-will in the powers of the mind, and there are powers
in Richard. The spirit will out, the dark spirit that
glows in his lust for dominion, his bitter feeling of hatred
and revenge; a volcano burns in him, and the volcano
must erupt. Atheistically, then, Richard
may in this respect be said to understand himself in evil
just as unconditionally as the ethicist understands himself in the
good: the directions are opposite, but the standpoint is
the same, the principle the same.

To maintain the principled difference between good
and evil, atheistic ethics must now seek to give the matter
another turn. It began by giving up
% --- p. 539 ---
\setcounter{footnote}{0}%
\opage{539}``the logical constraint of thought''; this rashness it must
bitterly regret. If it is to be able to master Richard, it must
begin, not personally with the subjective unity of
reason and nature, but impersonally with the objective unity
of reason and thinking. If it is true, as Fichte
says, that moral self-consciousness finds itself ``free from
all influences of the sensible world, absolutely active in itself,
thus as a power exalted above everything sensible,''
then it must give up existence, take refuge under the
``universal'' of metaphysics, and carry the struggle over into the logical
element of knowledge. The self of which the ethicist thinks, when
in the name of the good he speaks of knowing himself and willing
himself from pure necessity, is then not the individual
but the universal. Only the sensuous self, bound by drives, is
individual; the true spiritual self is universal,
it is a theory transparent in itself. Not what the individual
wills, but what reason wills through the individual, is the true
good; he who thinks the good cannot but will
it; to make oneself into the universal is intellectual love
of the good, only in such a love is there self-understanding.
Richard's fault is now this, that he is unwilling theoretically
to make his self into the universal; he loves
nothing, no one, in spite of his enormous self-love not even
himself: from that comes evil. His inflexible
defiance of the divine order of morality is
joined with a pain, a secret anguish, a restless
longing, which in the bitterness of his soul he does not understand,
because he does not wish to understand it. The three unhappy women,
Queen Margaret, the Duchess his mother, and
Queen Elizabeth, forget for a moment their mutual
grudge and unite as three Furies, in order to
% --- p. 540 ---
\setcounter{footnote}{0}%
\opage{540}pour out a stream of curses over Richard, who just
passes them by; but he will hear nothing. ``Sound
trumpets! --- Beat drums to the alarm! Let not heaven
hear those gossiping women who rail at the Lord's anointed;
strike, strike!'' As by defiance and the beating of drums he drowns out
the cry of the curse that sounds from without, so he seeks
by shrewd, sophistical reasoning to guard himself
against the curses from within and to drown out the Furies he
harbors in his inner being. With all the clarity of his penetrating
understanding Richard is not able to disperse his
moral darkness; his self-understanding is a continual self-stupefaction:
his deeper rational self does not come to
light.

Such soaring general considerations have something
flattering for the moral feeling; the only pity is that they
stand in conflict with the principle. In order to maintain the Idea
of the good, the atheistic ethicist will wrestle with Richard as
Hercules wrestled with the giant Antaeus, will lift him up
from the maternal earth of existence and strangle him in the airy region of
general thinking; but he overstrains himself
on the character; Richard is too heavy for the ethicist: it is
precisely the principle that weighs him down. In the abstractly general element
one can think, but not will, still less
act. The universal human being cannot act; for the universal human being
is a concept, and a concept cannot act.
In order to determine the real, the universal must go
through the individual; but when the universal goes
through the individual, then necessity enters
under freedom, thought under the will, and then
Richard has a secure footing. Even the reproach
that there is an inner voice he drowns out, a thought
he does not think because he does not wish to think it, he
% --- p. 541 ---
\setcounter{footnote}{0}%
\opage{541}puts to his own account; for by this it is without further ado admitted
that the thoughts of moral character are thoughts of the will,
that the course of thoughts depends on the will. With this admission
Richard is content; his thought is in all his
actions equal to itself and goes consistently as one with
his will; self-will is his law, the only law
he does not violate: to will in everything his own will is his
honor. But on these terms he is even in literal
agreement with the ethicist.

``There is something,'' says our ethicist, ``that counts for
me above everything, and for the sake of which I set everything else
aside, from whose assertion I let myself be deterred by no possible
consequence, something for which I would without
hesitation sacrifice my whole earthly welfare, my
good name and reputation, my life, the weal of the whole world,
if it could come into conflict with it. I will call
it honor.'' This Fichtean pledge Richard will
at once sign without demur; indeed, he will positively
be moved and fall in rapture into the ethicist's arms,
when he hears the closer explanation. ``This honor
I by no means place in the judgment of others on my actions,
not even if it could be the unanimous judgment of my age and
of posterity; but in the judgment I myself
pass upon them.'' That rings true! Richard listens attentively
and hears further: ``The judgment I pass upon my
actions depends on whether through them I remain in agreement
with myself, or through them come into
contradiction with myself. In the first case I can
approve them; in the latter case I would be dishonored
in my own eyes''\dots{} ``Consistency, agreement with
oneself!'' how it rings in Richard's ears! To
% --- p. 542 ---
\setcounter{footnote}{0}%
\opage{542}drown out such lovely voices he surely wants no
drummer. The universally rational, of which the atheistic
ethicist makes so much, can likewise cause
him no difficulties. My resolutions are ethical
when they are of such a kind ``that to the best of my
knowledge I can seriously will that all rational beings in
the same position should make the same resolutions.''
That is surely a lawful defense for Richard. Anyone
who thinks himself to be of the same nature as he, set in an age
filled with crimes, with his powers,
in his place, in his stead, must surely tell himself that
at bottom he would have to act as he does. When in a
civil society there is war of all against all, what counts
then but the right of the stronger? In this respect
Richard does only what reason commands: he strikes his
enemies to the ground, because he recognizes that each of
his enemies, if he were in his place, rationally
would have to strike him to the ground: he acts consistently.
The consequence of his way of acting by no means entails
the destruction of society; it leads only to this, that the stronger
defeats the weak, until the contending forces
come into equilibrium. Hypocrisy, slander, perjury and murder
are means that ought to be used only where they are needed; they
fall away of themselves when the times have become different.
How consistent Richard is can be seen from his relation
to Buckingham. This henchman of tyranny, who
until now has followed Richard as friend and counselor, suddenly
gets scruples and will not let himself be used to murder
the young princes; that is an inconsistency: by coming
into contradiction with Richard Buckingham comes into contradiction
with himself: no wonder that such a
% --- p. 543 ---
\setcounter{footnote}{0}%
\opage{543}contradiction costs him his head. Richard, by contrast, is,
as the principle demands, consistent in each and every respect, so
inflexibly consistent that no Stoic with his ``self-agreement''
as battle cry can dispute his rank. The hardness of his heart proves precisely that he,
exalted above sensuous compassion, feels himself only ``as
a spiritual power.'' He really approaches the theory
according to which the self's thinking is to be able, with the absolutely
objective necessity of knowledge, to pass over into will, into
resolution, into action. He acts not from mere natural impulse
or from sheer wickedness, but according to a strictly
thought-through and deeply laid plan; in acting with
freedom he does only what is necessary. Henry
the Sixth is murdered, it is necessary; Clarence is murdered,
it is necessary; Edward's sons are murdered, it is altogether
necessary; Buckingham's head falls, that too
is necessary. And what then is freedom? Precisely what
an atheistic metaphysics so ingeniously explains to us, the developed
concept of necessity, its immanent transition
into rational self-consciousness! And so a hero
like Richard, a character who satisfies the
theory to so high a degree, is not to be said to understand himself,
even if his self-understanding is in evil?

For the atheistic ethicist there is then only the way out
left of giving up the merely formal, laying the emphasis
on the content of morality, and appealing to the Idea.
What agrees with the Idea is good, what conflicts
with it is evil; it is not the consistent way of acting
in and for itself, but the character's relation to the Idea
as Idea, that in ethics is decisive. Richard's blindness
then shows itself in this, that he convulsively clings to
% --- p. 544 ---
\setcounter{footnote}{0}%
\opage{544}his egoistic self and violates the Idea in every shape;
friendship, family, state, everything he sacrifices to his
egoism. He is worse than a Machiavellian; for
the Machiavellian does still posit the Idea in the form of the state as
the absolute, but King Richard pays no heed to the general
weal, he takes the state as a prey: for him nothing is
absolute but his own domineering ego. This break
with all moral powers is the deep contradiction in
Richard's being, and the consequence of this break a
breakthrough that ends in despair. After his accession to the throne,
even before rumor tells him that Henry
of Richmond has landed with an army, Richard's self-confidence
begins to decline. His strength is
broken, his soul has no peace, he is sleepless; sleep
itself, which otherwise refreshes the weary, has, when
it finally comes, no refreshment for him, it
torments him with heavy dreams; gnawing memories
in streams of blood, voices from the graves crying: Revenge!
The rending of the brotherly bond, the mother's curses,
full certainty of everyone's enmity, everyone's loathing, everyone's
hatred, all this together oppresses Richard's mind, cows his
courage, and lets him feel that there is a moral power that
punishes the evil one. Thus the last catastrophe approaches.
Richard sleeps in his tent on ``Bosworth
Field'' and dreams. The spirits of the murdered appear to
him, one after another, each weighs upon him with his
memory, his lament, and through all memories, all laments
rings the word of judgment: ``despair and die!'' Terrified Richard
starts up,---``it is dead midnight; the cold drops
bubble on his body.'' Then follows the famous monologue,
in which for the first and also for the last time he
% --- p. 545 ---
\setcounter{footnote}{0}%
\opage{545}gets his ethical self to speak. Now at last he understands
himself, but the self-understanding is the most dreadful
self-contradiction. ``What do I fear? Myself? I am
alone: yes Richard loves Richard---I am I. Is there
a murderer here? No;---yes I am one, then flee!---What, from
myself? Good ground: why? I will avenge. Whom?
Myself, upon myself? I love indeed my own self; why?
for all the good I have done myself. O, no!
Alas no! I ought rather to hate myself for hateful deeds against
my self.'' The ethical self-contradiction is a literal
self-condemnation; to get Richard condemned we need not
appeal to a supramundane right; the court is already in session:
it is no omniscient God on the Day of Judgment, but
simply and plainly the moral Idea that condemns Richard.
Surrounded by the terrors of judgment he stands as transformed.
When he began by declaring that he would
be a villain, it was only a presumptuous mockery; there
was in his consciousness no moral difference between a
villain and a noble hero; but now, when he is struck by
the lightning of the spirit and in his soul's anguish cries out: ``I am a
scoundrel!'' now the difference between evil and the
good has dawned on him: the Idea has come to revelation,
but in order to crush him. How crushed he
is we hear from his despair, his soul-shaking
self-confession; for ``his conscience has a thousand tongues,
each tongue brings in a separate tale, each
tale condemns him as a villain. Perjury and most false
oath in the highest degree, murder, cruel murder in the most fearful
degree, and all sins in every degree, to the bar
crowd and cry: Guilty, Guilty!''
If it were now really the case that the
% --- p. 546 ---
\setcounter{footnote}{0}%
\opage{546}Idea of the good alone by itself, without ground in God's will,
has got Richard condemned and done justice its due,
then atheistic ethics has by this immanent
act of judgment undeniably celebrated a great triumph. It is then
made known once for all that the difference between
good and evil can not only subsist, but even first
become objectively grounded, when God is denied, and
the unclear notion of a judge beyond is completely
given up. But if in any scientific
inquiry there is ground to speak of ``smuggling in,'' ``eclecticism,''
``self-contradiction'' and the like, it is surely here.
The whole grounding rests on a smuggling in, in that
one imperceptibly lets the \emph{moral}, immediately moral
\emph{feeling} take the place of the \emph{moral principle.}
That the moral feeling is wounded when evil is put on a par
with the good is certain, just as it is also certain that
the poet has poetic right in letting Richard's better
nature break through and break his defiance.
But that is surely not at all what is at issue here. The question,
the question the atheist here ought to answer,
concerns not the feeling but the principle, the principle by
virtue of which the opposition between good and evil
alone has validity. Is Richard really condemned by virtue
of the atheistic principle? By no means! He is in the proper
sense not morally condemned at all; he is merely frightened,
overwhelmed by dream visions; what he says, he says, if
not half asleep, still in a state in which he is not
master of himself and his principle. As soon as he is properly
awake and has come to his senses, he also disavows
all that he said in the first fright. ``Conscience,
you coward, you terrify!'' This
% --- p. 547 ---
\setcounter{footnote}{0}%
\opage{547}opening line indicates clearly enough where the following stream of
self-reproaches will lead. Richard has had a weak
moment, and weak moments even the strongest can have,
whether he is otherwise good or evil. To Ratcliff, who
has just entered the tent, he says that he has had
a dreadful dream; he ``fears, fears''! To Ratcliff's
encouraging words: ``My good prince, be not afraid of
shadows,'' he answers: ``Now by the apostle Paul, the shadows
of the night have struck more terror to Richard's soul
than ten thousand real warriors, armed,
shot-proof, under Richmond could.'' What kind of
terror it is with which he is struck, and what kind of
fear he means, is clear enough: it is by no means fear
of the judgment of conscience, but only fear of losing
the battle. Had Richard been morally overcome in his principle,
he would surely have had to be brought, in full possession of his senses,
to give the Idea its due and to acknowledge conscience. But
far from it. Scarcely has he risen, got his armor
on and is about to begin the battle, before we hear in his voice of command
once more the wild defiance: ``Go, my men,
each to his business, let no gossiping dreams
frighten the soul; conscience is a word cowards use,
introduced to curb the strong. Let our strong arm be
conscience, sword our law; so charge forward!'' So
the ``Idea'' shall not boast that it
has overcome Richard in the name of conscience!

``But,'' it is said, ``precisely because the Idea cannot conquer
Richard subjectively, it works as an objective power and
conquers him objectively.'' Now the ethical has been left behind,
and the misunderstanding is in full swing; for an arrow, a bullet,
a stone, a flash of lightning can also conquer an individual
% --- p. 548 ---
\setcounter{footnote}{0}%
\opage{548}objectively: what an honor for the Idea to be set on an equal
footing with a mere force of nature! The Idea, it is said, is
not a force of nature, it is a force of spirit, spirit itself in history,
in the world of freedom; as spirit present in and immanent to individuals,
actions and events it governs
everything and everyone. It is only in his own imagination that
Richard has been able to tear himself loose from the Idea; he is in its
power; it holds him fast with an invisible hand from first
to last, it uses him as an instrument. Among the
powerful, who in that period of terrors had all offended
against society, Richard, the blood-dripping
hero who hews off the heads of the hydra of civil war,
is precisely by the power of the Idea the most powerful. ``He is,'' says
Rötscher, ``as the Nemesis of the World-Spirit. By this view
Richard already gains an exalted pedestal.
His foundation is the torn-up state, his goal is the crown;
his absolute significance is to be the bearer of
the historical spirit, by whom all those heads are mown off
which have undermined the universal order, dissolved the state and
burdened themselves with crimes. In leveling everyone's
right, he levels the ground, so that the unity of the state may once more
raise itself out of the discord of interests\dots{} But
as a product of the whole, dissolving state Richard
himself perishes, for the crime must be atoned for and
the tool cast away, as soon as it has accomplished its
bloody work.'' From a historical-aesthetic standpoint
there is nothing to object to in this judgment; but
what becomes of ethics? Although Richard
is just not on any objective familiar footing with
the Idea as Idea---such an objective, intimate familiarity
must be reserved for the armchair metaphysician---yet he is
% --- p. 549 ---
\setcounter{footnote}{0}%
\opage{549}for his time the Idea's most zealous servant, its most faithful
henchman, the most untiring worker in its vineyard;
understood objectively, the spirit is also far
more content with him than even with the pious King
Henry, not to speak of the others. Had Richard
been a saint or at least a noble, upright man,
the spirit of history could surely not have used him
for its bloody work. But if Richard has done only
what the Idea demanded: by what law can the Idea then
condemn him? It does happen sometimes that a ruler
casts away the instrument that thought it deserved
thanks for having fulfilled its lord's wish. Thus
Bolingbroke (Richard II, Act 5, end) says to Exton,
who brings him the body of Richard the Second in the coffin:
``Expect no thanks from me! for you have laid with bloody
hand a shame without equal upon me and my renowned
realm.'' Exton reminds the King of his unequivocal
exhortation: ``Your word, my lord, occasioned my deed,''
to which Bolingbroke (Henry IV) replies: ``He loves
not poison who uses poison; I do not love you; his death was
needful to me; but you I hate, love him in death;
the sting of conscience shall be your wage, but no praise and
princely favor from me; go wander friendless like Cain in the horror
of the night, and show not your head to the light of day!'' Noble feeling,
high pathos, a splendid language; poetry exults, but
---ethics weeps.

If the Idea really could treat Richard as
Bolingbroke treated Exton, things would look black
for the distinction between good and evil. Who is it that bears
the guilt, ``the poison'' or ``he who uses the poison?''---``The crime
must be atoned for!'' but surely not the crime of
% --- p. 550 ---
\setcounter{footnote}{0}%
\opage{550}having satisfied the Idea? When on the atheistic side
Richard's crimes are spoken of, one speaks plainly
and simply subjectively out of natural human feeling, but
by no means metaphysically objectively out of the ``spirit of history.''
It is on the whole characteristic of godless
theories in all their forms that, after having put
the universally rational in God's place, one is compelled to personify
the Idea and, in order to preserve a semblance of morality,
to ascribe to it a kind of conscience and to treat it as
a moral personality. One insists on the concept and smuggles in
notions that stand in jarring conflict with the concept;
one deifies the abstraction and paints up the dead
abstract to look like a moral being that
rewards the good and punishes the evil: it is a mystification
equally ghastly for life and for science. The Idea, the Idea torn loose from
God's holy, righteous will, is no
moral power, neither good nor evil, neither just
nor unjust. Its essence is in the order of nature, in
the working of forces, in the necessity of grounds: it lives
in the fulfillment of objective laws. Its right is a spiritual
natural right: everything that is possible is permitted, only the
impossible is forbidden. Evil too happens according to laws;
therefore the Idea is just as reconciled with evil as with
the good; it is without joy, but also without pain,
without wrath, but also without compassion; it is not indignant
at violence and plunder, it does not grieve when an executioner
is crowned, any more than the ax sighs because an innocent man
is executed.

The result of our inquiry is this: from the standpoint of atheism
there is just as clear and thoroughgoing a
self-understanding in evil as in the good, and
% --- p. 551 ---
\setcounter{footnote}{0}%
\opage{551}that for the simple reason that the atheistic principle is not able
to maintain the moral difference between
good and evil. It is no use appealing to
the nature of man, for human nature exists only in real
human beings, and real human beings each have their own
nature, in so far as they each have their own mind. He who has
a good mind prefers the good, he who has an evil
mind, evil; if morality is to rest immediately
on the mind, then the evil man is just as justified as the
good. It is also no use appealing to man's
rational being, for the thought that is to grasp the
universally rational must,---if it does not wish, through empty
theories, to exclude itself from life and reality,---
of necessity place the rational being in man's
thinking and willing self and make consistent
self-satisfaction the supreme principle. But when
self-will is law and consistency the measure, the
evil man can in his way have just as consistent a self-satisfaction
in carrying through his will as the good man in
carrying through his. The satisfaction is on neither side
perfect; the evil man has a resistance in his better nature
and thereby in himself; likewise the good man has a resistance in his
lower nature and thereby in himself; the condition for the one as
for the other is unavoidable struggle and strife. How great the satisfaction
will be depends solely on the strength of character
with which the principle is carried through. The consistent character
has a great, the inconsistent a slight self-satisfaction, without any regard to good and evil. After
having made the principle as subjective as possible and placed
the universal in self-will, it is in vain to want to bring
the matter before an objective forum and let the Idea judge.
% --- p. 552 ---
\setcounter{footnote}{0}%
\opage{552}The moral Idea is an Idea of the will; if it does not have its center
in God's will, then it is unconditionally dependent on
man's will; the Idea's judgment on man is then
only man's own judgment on himself: the judgment is
unconditionally subjective. If the Idea is nonetheless to raise
itself above the subjective and count as a self-valid objective
power, then it ceases to be ethical.

Hereby it is by no means proved that an atheist
could not be a moral character, quite the contrary! it
is precisely proved that, according as his mind and subjectivity
are, he can be both moral and immoral, since it
is shown that the opposition of the moral and the immoral is, not
for the individual atheist, but for the atheistic principle,
without significance.

There are, as is well known, in art, poetry and history
certain great orienting points of view of which the
cultivated man, however little he may otherwise have occasion to immerse
himself in any of the said subjects, will not easily forgive
himself for being wholly ignorant. Ought the
same not also to hold when the talk is of religion
and morality? Among the principal propositions
orienting in this respect we count the following:

God's existence, the right of conscience and the soul's
immortality are personal presuppositions, belonging so
unconditionally to the inner life of the soul that their validity, when
we take it strictly, can no more be proved than disproved
by science. They are propositions of faith, whose
peculiarity science may well illuminate, but not
ground.

Theology, which wishes to ground, and atheism, which
wishes to undermine revealed religion, are both equally
% --- p. 553 ---
\setcounter{footnote}{0}%
\opage{553}much in the wrong. The unmistakable superiority of atheism
rests solely on this, that the struggle is conducted in the element of knowledge,
an element that stifles faith. The impotence of theology is
no proof against religion, for religion with its principle
has nothing to do with theology.

The true personal proof of the validity of morality
is inseparable from the true personal proof of
the validity of religion; the proof rests on the unconditional
reality of the will. The human will loses
instantly its true stamp of eternity, its absolute dependence
and with it also its higher moral independence,
when faith in God's will is excluded.

Only when the religious-moral principle is recognized
in complete independence of objective knowledge and
science, can it with all-sided consistency unfold the content
of its cognitions in its own sphere, and the doctrine of religion
arise reborn to renewed activity. To wish, from
an atheistic standpoint, to attack this view by
lashing out with ``scientific contradictions'' is
nonsense: since the element of knowledge, according to our chosen standpoint,
is opposed to the essence of religion, it follows
quite of itself that the doctrine of religion, when it is treated
in science, must necessarily swarm with contradictions.
The only way to be rid of the contradictions is
to separate religion from science.{\renewcommand{\thefootnote}{*)}\footnote[1]{Cf. Philosophy of Religion pp. 5—11, 74---81, 97—108 ff.}}

By holding fast to these principles cultivated
readers will easily be able to orient themselves in the confusion of theological
and atheistic assurances that are put into circulation
with an appeal to ``grounds,'' ``clear grounds,''
% --- p. 554 ---
\setcounter{footnote}{0}%
\opage{554}``irrefutable grounds.'' It is one thing to cry out with
grounds, another to give grounds; the former is easier,
but the latter is more reliable. Those who cry out most eagerly
with grounds are not always reliable, not so reliable
as those who cry out with herrings: the latter
have at least as a rule the ware they offer for sale, but the former
unfortunately often have nothing but the cry.

\signature{R. \emph{Nielsen}.}

\end{document}
